\chapter{机械臂控制}
\label{ch:arm-ctrl}

% ════════════════════════════════════════════════════════════════
%  本部第 4 章：机械臂控制。定位（最重要纪律）：
%  P4「面向机器人的控制」(sec:ctrl-robot) 与 P4 操作空间章 (ch:operational-space)、
%  力控章 (ch:force-wbc) 已覆盖：动力学方程/计算力矩(带定理)/PD+g 无源性/OSC/阻抗/WBC。
%  本章 = (1) 凡基础定理一律 \cref 回 P4 不重证；(2) 深化 P4 只简述处（OSC 在机械臂上的
%  完整推导、阻抗整定、力位混合的选择矩阵理论）；(3) 补 P4 没有的「机械臂专精」
%  （异质惯量臂的控制率选择、增益异质设计、力矩臂真阻抗范式、病态质量阵对 OSC 的
%  数值影响）；(4) 用六轴力矩臂 arm_dm 的实测对比作实证血肉。
%  理论网络重建（本部经典理论无教材抽取源）：khatib1987operational / hogan1985impedance /
%  takegaki1981feedback / raibert1981hybrid，已网核 OSC 动力学一致逆/Hogan 阻抗/选择矩阵公式。
%  实践落到 \cref{ch:arm-prac}（已写好）。复用第 1/2 章雅可比/条件数 \cref{ch:arm-kin,ch:arm-dyn}。
% ════════════════════════════════════════════════════════════════

\begin{practice}[前置自测]
开始之前，先试答下列问题。若已能流畅作答，可只速读本章的对比表与速查卡；若卡住，正是本章要为你解决的。
\begin{enumerate}
  \item 上一部分（\cref{ch:control}）已经证明：对\emph{点到点定位}，只需重力补偿 $\mathbf{g}(\mathbf{q})$ 加一组 PD 就能全局稳定，连质量矩阵 $\mathbf{M}$ 都不必精确知道。可一旦任务从「停在一点」变成「跟住一条时变轨迹 $\mathbf{q}_d(t)$」，为什么 PD+重力补偿就会出现\emph{周期性滞后}？要消掉这个滞后，控制律里必须补进哪两项？
  \item 你把一台六轴臂的关节增益统一设成 $\mathbf{K}_p=150\mathbf{I}$，肩肘跟得很好，可手腕（真实惯量仅 $2\times10^{-4}$）却高频发散。诊断发现：对这只腕施加的力矩在钳位前竟冲到 $888\,\mathrm{N\!\cdot\!m}$。一个这么\emph{轻}的关节为什么\emph{更难}控、而不是更好控？为什么计算力矩里的 $\mathbf{M}(\mathbf{q})$ 加权能天然救它，而均匀增益 PD 不能？
  \item 操作空间控制（OSC）让你「在末端坐标里直接施力」、同时把次任务藏进零空间。教科书给的零空间投影是 $\mathbf{N}=\mathbf{I}-\mathbf{J}^\top\bar{\mathbf{J}}^\top$，里头那个 $\bar{\mathbf{J}}$ 偏偏含 $\mathbf{M}^{-1}$。如果你图省事把它换成纯运动学的 Moore--Penrose 伪逆 $\mathbf{J}^+$，末端会被零空间力矩\emph{污染}出一个寄生加速度——这是为什么？这条「动力学一致」的逆，凭什么是\emph{唯一}让零空间力矩不扰动末端的逆？
  \item 一台\emph{力矩控制}的臂（电机底层只收力矩指令）想让末端像弹簧一样：被推就退、撤手就弹，且\emph{不装任何力/力矩传感器}。这件事物理上为什么可能？阻抗控制 $\boldsymbol{\tau}=\mathbf{J}^\top(-\mathbf{K}_c\Delta\mathbf{x}-\mathbf{D}_c\dot{\mathbf{x}})+\mathbf{g}$ 里，把 $\mathbf{K}_c$ 调小到底改变了末端的什么\emph{物理量}？为什么阻尼 $\mathbf{D}_c$ 取 $2\sqrt{\mathbf{K}_c}$、而真要在每个笛卡尔方向都临界阻尼时这个公式还不够？
  \item 打磨、装配这类任务要「某个方向控力、另一个方向控位」。Raibert--Craig 的选择矩阵 $\mathbf{S}$ 是怎么把一个六维空间\emph{劈成}「力子空间」与「位子空间」两半的？为什么这种「劈分」必须在一个精心选定的\emph{约束坐标系}里做，而不能随便沿世界系的 $x,y,z$ 劈？
  \item 你的力矩跟踪坏了，示波器上关节波形有四种典型长相：卡在初始位置一动不动 / 朝一个方向越跑越快（反号失控）/ 高频抖到发散 / 跟得上但总慢半拍。不看代码，光凭这四种波形，你能各自反推出大概是哪一类根因吗？
\end{enumerate}
\end{practice}

\section*{本章目标}
学完本章，你应当能够：
\begin{enumerate}
  \item \textbf{接续}上一部分的控制工具（\cref{sec:ctrl-robot}）：说清「关节空间 vs 任务空间」「调节 vs 跟踪」「无模型鲁棒 vs 有模型精确」三组分工，并画出本章相对 P4 的\emph{增量地图}（\cref{sec:actrl-bridge}）；
  \item \textbf{把} PD+重力补偿与计算力矩\emph{落到机械臂}：复用 P4 已证的无源性/线性化定理（不重证），补足「跟踪场景」的误差动态与实测——计算力矩比 PD+重力补偿在跟踪上紧 $26\times$，且看懂这个差距正是前馈 $\mathbf{M}\ddot{\mathbf{q}}+\mathbf{C}\dot{\mathbf{q}}$ 的功劳（\cref{sec:actrl-pdg,sec:actrl-ctc}）；
  \item \textbf{诊断并设计}异质惯量臂（重肩肘 + 轻腕）的控制率：用「自然时间常数 vs 控制频率」「质量阵加权给力矩权限」两条第一性原理，解释为何低惯量腕必须\emph{解耦}、为何增益要\emph{异质}，并算出腕误差如何经质量阵耦合腐蚀主关节（\cref{sec:actrl-heterogeneous}）；
  \item \textbf{完整推导}操作空间控制（OSC）：从任务空间动力学 $\boldsymbol{\Lambda}\ddot{\mathbf{x}}+\boldsymbol{\mu}+\mathbf{p}=\mathbf{F}$、对偶映射 $\boldsymbol{\tau}=\mathbf{J}^\top\mathbf{F}$，到动力学一致广义逆 $\bar{\mathbf{J}}$ 与零空间投影 $\mathbf{N}=\mathbf{I}-\mathbf{J}^\top\bar{\mathbf{J}}^\top$，并理解它对质量阵\emph{条件数}的高度敏感（\cref{sec:actrl-osc}）；
  \item \textbf{整定}阻抗与导纳控制：区分「测位置出力」与「测力出位置」，推导笛卡尔阻抗律、临界阻尼的设计、以及「力矩臂 $\to$ 真阻抗无需 F/T」的范式，并用 $50\,\mathrm{N}\to0.2\,\mathrm{m}=F/K$ 的线性弹簧律实证刚度（\cref{sec:actrl-impedance}）；
  \item \textbf{构造}力位混合控制：用选择矩阵 $\mathbf{S}$ 在约束坐标系里把方向劈成力/位两半，理解自然约束与人工约束的对偶、以及为何力矩臂上要用\emph{全雅可比}而非关节切分（\cref{sec:actrl-hybrid}）；
  \item \textbf{运用}一套力矩跟踪的\emph{故障诊断签名}：把波形归类为 stuck-at-home / wrong-sign-runaway / unstable / lagging 四类，并映射到根因（\cref{sec:actrl-diagnostics}）；下接 \cref{ch:arm-prac} 的全栈实战与 \cref{ch:rlmc-manip} 的强化学习操作。
\end{enumerate}

\subsection*{本章知识导航}
本章是一条「\emph{从关节空间调节器，逐步武装到任务空间力控、再武装到接触}」的主线。它\emph{不}重证上一部分已立的基础定理（计算力矩线性化、PD+重力补偿无源性都在 \cref{sec:ctrl-robot} 证过），而是把那些「一张图景」级别的结论，在机械臂这个具体载体上\emph{展开成完整工程谱系}：补全 OSC 的逐步推导、阻抗的整定、力位混合的选择矩阵理论，并补进 P4 完全没有的「异质惯量臂控制率选择」这一工程专精。结构如下：

\begin{center}
\begin{forest} navtreeLR
[{机械臂控制\\（承 P4 深化 + 专精）}
  [{关节空间\\调节与跟踪}
    [{PD+g / 计算力矩\\CT vs PD+g 26$\times$}] ]
  [{异质惯量臂\\控制率选择 $\bigstar$}
    [{增益异质 / 解耦轻腕\\$M$ 耦合腐蚀}] ]
  [{任务空间\\OSC 完整推导}
    [{动力学一致逆\\对 cond($M$) 敏感}] ]
  [{接触控制\\阻抗 / 导纳 / 力位混合}
    [{力矩臂真阻抗\\选择矩阵 $S$}] ]
  [{故障诊断\\四签名映射}
    [{$\to$ 实战 / RL 操作}] ]
]
\end{forest}
\end{center}

\subsection*{前置桥接}
本章把工具借自三处、并把成果交给两处，列清依赖以便读者回查：

\begin{itemize}
  \item \textbf{承上（控制理论地基）}：机器人动力学方程与结构性质 \cref{eq:ctrl-manip}、\cref{prop:ctrl-robot-props}（$\mathbf{M}\succ0$、$\dot{\mathbf{M}}-2\mathbf{C}$ 斜对称）；计算力矩与其线性化定理 \cref{sec:ctrl-robot-ctc}、\cref{thm:ctrl-ctc}；PD+重力补偿的无源性证明 \cref{sec:ctrl-robot-pd}、\cref{thm:ctrl-pd-grav}；操作空间与阻抗的入门图景 \cref{sec:ctrl-robot-osc}、\cref{sec:ctrl-robot-impedance}、\cref{def:ctrl-impedance}。Lyapunov/LaSalle 工具见 \cref{ch:lyapunov-basics}。
  \item \textbf{承上（运动学/动力学）}：雅可比（几何/解析/旋量）、静力学对偶 $\boldsymbol{\tau}=\mathbf{J}^\top\mathbf{F}$、奇异与阻尼最小二乘见本部 \cref{ch:arm-kin}；质量矩阵 $\mathbf{M}(\mathbf{q})$、操作空间惯量、\emph{质量阵条件数}的第一性原理见 \cref{ch:arm-dyn}（尤其 \cref{sec:ad-conditioning}）；位姿、$\SE(3)$、右扰动/对数映射见 \cref{ch:rigid_body,ch:lie}。
  \item \textbf{平行（不重复，互链）}：P4 已有专章把 OSC（\cref{ch:operational-space}）与力控/WBC（\cref{ch:force-wbc}）展开到浮动基、TSID、稳定性证明；本章只在\emph{固定基串联臂}视角下给「自包含的完整推导 + 机械臂专精 + 实测」，与那两章互为补充而非重复。
  \item \textbf{启下}：本章控制律落到真工程见 \cref{ch:arm-prac}；经典控制如何支撑强化学习操作见 \cref{ch:rlmc-manip}。
\end{itemize}

\noindent 全章记号统一为：粗体 $\mathbf{q},\boldsymbol{\tau}\in\mathbb{R}^n$ 为关节角与关节力矩，$\mathbf{x}\in\mathbb{R}^m$ 为末端任务坐标，动力学方程
\begin{equation}
  \mathbf{M}(\mathbf{q})\ddot{\mathbf{q}}+\mathbf{C}(\mathbf{q},\dot{\mathbf{q}})\dot{\mathbf{q}}+\mathbf{g}(\mathbf{q})=\boldsymbol{\tau}+\boldsymbol{\tau}_{\rm ext},\qquad \boldsymbol{\tau}_{\rm ext}=\mathbf{J}^\top\mathcal{F}_{\rm ext},
  \label{eq:actrl-manip}
\end{equation}
与 \cref{eq:ctrl-manip} 完全一致（$\boldsymbol{\tau}_{\rm ext}$ 为外接触力的广义力）。

% ════════════════════════════════════════════════════════════════
\section{控制问题与对 P4 的衔接}
\label{sec:actrl-bridge}

\paragraph{动机：机械臂控制不是「再学一遍 PID」。}
读到这里，读者已在 \cref{ch:control} 见过反馈、LQR、MPC、Lyapunov，也在 \cref{sec:ctrl-robot} 见过机器人动力学方程与两大范式。一个自然的疑问是：既然那一节已经把计算力矩、PD+重力补偿都证过了，机械臂控制还剩什么可讲？答案是：上一部分给的是「\emph{一张图景}」——告诉你这些控制律\emph{存在}且\emph{为何稳定}；本章要给的是「\emph{一套工程}」——告诉你在一台\emph{具体}的、惯量分布\emph{不均匀}的真臂上，这些律\emph{怎样落地、何处会坏、坏了怎么救}。两者的关系，恰如「会证明牛顿第二定律」与「会调一台真发动机」。

\paragraph{两个坐标空间。}
机械臂控制的全部问题，都可归结到「在哪个空间里写控制目标」：

\begin{definition}[关节空间控制与任务空间控制]
\label{def:actrl-spaces}
\textbf{关节空间控制}把目标写成关节量 $\mathbf{q}_d(t)$，直接对每个关节施控；末端位姿由正运动学被动决定。\textbf{任务空间（操作空间）控制}把目标写成末端量 $\mathbf{x}_d(t)$（笛卡尔位姿/力），控制律先在任务空间算出所需广义力 $\mathbf{F}$，再经对偶映射 $\boldsymbol{\tau}=\mathbf{J}^\top\mathbf{F}$ 落到关节力矩。
\end{definition}

\begin{insight}
两个空间的取舍是机械臂控制的第一性决策。\textbf{关节空间}：直接、解耦、无奇异顾虑、适合「按既定轨迹走」（轨迹已由 \cref{ch:arm-traj} 规划好，控制只需跟住关节路点）；缺点是末端误差被运动学非线性放大、不便表达「末端要顶住多大力」。\textbf{任务空间}：目标即末端行为（位置/力/阻抗），是接触任务（装配、打磨、柔顺）的唯一自然语言；代价是要算雅可比、近奇异要稳健化、且——本章的核心教训——\emph{对质量矩阵的数值条件极其敏感}。一条实用经验：\emph{自由空间快速移动用关节空间（计算力矩），接触/柔顺用任务空间（OSC/阻抗）}；全栈系统两者接力（见 \cref{ch:arm-prac} 的 capstone）。
\end{insight}

\paragraph{本章相对 P4 的增量地图。}
为避免与上一部分重复、又让读者一眼看清「新增了什么」，\cref{tab:actrl-delta} 把本章每节相对 P4 的关系标清：凡 P4 已证的，本章只交叉引用回去取用；凡 P4 只简述的，本章给完整推导；凡 P4 没有的，本章原创补足。

\begin{table}[H]\centering\small
\caption{本章相对 P4「面向机器人的控制」的增量}
\label{tab:actrl-delta}
\begin{tabular}{p{0.22\textwidth} p{0.30\textwidth} p{0.40\textwidth}}
\toprule
本章节 & P4 已有（交叉引用取用） & 本章增量 \\
\midrule
\cref{sec:actrl-pdg} PD+g & \cref{thm:ctrl-pd-grav} 无源性证明 & 落到臂、重力补偿验通路、作 CT 的跟踪基线 \\
\cref{sec:actrl-ctc} 计算力矩 & \cref{thm:ctrl-ctc} 线性化定理 & 跟踪场景误差动态、$26\times$ 实测、公平对比设计 \\
\cref{sec:actrl-heterogeneous} 异质惯量臂 & ——（P4 \emph{无}） & \textbf{专精}：控制率选择、增益异质、解耦轻腕、$M$ 耦合腐蚀 \\
\cref{sec:actrl-osc} OSC & \cref{sec:ctrl-robot-osc} 简述 & \textbf{完整推导}：动力学一致逆 + 零空间 + 对 cond($M$) 敏感 \\
\cref{sec:actrl-impedance} 阻抗/导纳 & \cref{def:ctrl-impedance} 定义 & \textbf{深化整定}：临界阻尼设计、力矩臂真阻抗、$F/K$ 实证 \\
\cref{sec:actrl-hybrid} 力位混合 & \cref{ch:force-wbc} 谱系（互链）& 选择矩阵理论 + 约束坐标系 + 全雅可比工程 \\
\cref{sec:actrl-diagnostics} 故障诊断 & ——（P4 \emph{无}）& \textbf{专精}：四签名 $\to$ 根因映射 \\
\bottomrule
\end{tabular}
\end{table}

% ════════════════════════════════════════════════════════════════
\section{关节空间：PD + 重力补偿}
\label{sec:actrl-pdg}

\paragraph{先问一句：凭什么能「一个关节一个关节地」控？}
本节与下一节都把关节当作\emph{一个个独立的回路}来设计——这件事并非天经地义。动力学 \cref{eq:actrl-manip} 里的 $\mathbf{M}(\mathbf{q})$ 是\emph{满阵}且\emph{随构型剧变}：关节 $i$ 的运动经非对角元 $M_{ij}(\mathbf{q})$ 推着关节 $j$ 动、对角元 $M_{ii}(\mathbf{q})$ 又随位形变好几倍。严格说，机械臂是个强耦合、强非线性的多体系统，「各关节独立」似乎根本不成立。可工业现场偏偏大量用「每关节一组 PD」就把活干得很好。这中间缺的，正是一条\emph{判据}：到底\emph{何时}可以把这台臂近似成一串解耦的单输入单输出（SISO）二阶系统、何时不能。补上它，本节的 PD（利用解耦）与下节的计算力矩（在解耦失效时显式补偿耦合）才各得其所。

\begin{insight}[独立关节（解耦 SISO）假设何时成立：高减速比把耦合按 $1/k_r^2$ 压平]
\label{ins:actrl-decouple-valid}
所谓\textbf{独立关节控制}，是把第 $j$ 关节近似成一个\emph{定常惯量、与其他关节无耦合}的二阶系统 $J_{\rm eff,j}\ddot q_j+b_j\dot q_j=\tau_j$，再各自闭一个 SISO 的 PD/PID。它\emph{何时有效}？关键在\emph{减速比} $k_r$。\cref{ch:drive-systems} 已逐步推出（\cref{eq:ds-Ieq-motor,ins:ds-hide-coupling}）：从\emph{电机侧}看，关节 $j$ 的有效惯量是 $J_{m}+M_{jj}(\mathbf{q})/k_r^2$——常数转子惯量 $J_m$ 占大头，而\emph{构型相关}的负载惯量被 $1/k_r^2$ 压成小扰动；同理关节 $i$ 对 $j$ 的耦合 $M_{ij}(\mathbf{q})$ 也被同一个 $1/k_r^2$ 压平。工业臂 $k_r$ 动辄数十到数百，$1/k_r^2$ 是 $10^{-3}$--$10^{-4}$ 量级，于是\emph{构型相关性与关节间耦合双双退化成可当扰动吸收的小项}，每个关节近似定常解耦——这就是「为何工业臂用独立关节 PD（甚至不必精确建模 $\mathbf{M},\mathbf{C}$）就够用」的理论判据（Craig\cite{craig2005introduction} §9.9、Spong\cite{spong2006robot} §8 的单关节模型正建立于此）。\textbf{何时失效}：当 $1/k_r^2$ 不再小到能压住耦合——\emph{直驱}（$k_r=1$）、\emph{低减速比}、或\emph{高速}（科氏/离心项 $\mathbf{C}\dot{\mathbf{q}}$ 随速度平方增大、不再可忽略）——构型相关惯量与耦合\emph{全部抬头}，独立关节假设破产，必须改用\emph{显式补偿动力学}的\textbf{计算力矩}（\cref{sec:actrl-ctc}）：它直接算进满阵 $\mathbf{M}(\mathbf{q})$ 与 $\mathbf{C}(\mathbf{q},\dot{\mathbf{q}})$，把耦合\emph{主动抵消}成解耦的误差动态。一句话定位本节与下节的分工——\textbf{独立关节 PD 是「靠齿轮箱替你解耦」，计算力矩是「靠模型替你解耦」}；前者是高减速工业臂的够用起点，后者是直驱/高速/高性能臂的必需。
\end{insight}

\paragraph{动机：最省的跟踪基线。}
机械臂控制最简单的有效律，就是「重力补偿 + 比例微分」。它\emph{不需要}质量矩阵、不需要科氏阵，只要重力 $\mathbf{g}(\mathbf{q})$——这恰是 \cref{sec:ctrl-robot-pd} 证过的：对点到点定位，\cref{eq:ctrl-pd-grav} 的 $\boldsymbol{\tau}=\mathbf{g}(\mathbf{q})+\mathbf{K}_p\tilde{\mathbf{q}}-\mathbf{K}_d\dot{\mathbf{q}}$ 凭无源性 + LaSalle 即全局渐近稳定，且对 $\mathbf{M},\mathbf{C}$ 的误差鲁棒。本章不重复那个证明，只补两件 P4 没做的事：把它\emph{落到一台真臂}，并把它当作下一节计算力矩的\emph{公平对照基线}。

\paragraph{重力补偿：验通「力矩权限」。}
在任何力矩控制实现里，第一个该跑的不是跟踪、而是\emph{纯重力补偿} $\boldsymbol{\tau}=\mathbf{g}(\mathbf{q})$（即 \cref{eq:ctrl-pd-grav} 取 $\mathbf{K}_p=\mathbf{K}_d=\mathbf{0}$）。它的诊断价值在于：若 $\mathbf{g}(\mathbf{q})$ 算对、力矩通路也通，机械臂应当\emph{失重悬浮}——在任意构型撒手不掉。六轴力矩臂 arm\_dm 的实测里，这一步确认了「强力矩权限」：仅 $8\,\mathrm{N\!\cdot\!m}$ 的补偿力矩就能驱动腰关节、$-6\,\mathrm{N\!\cdot\!m}$ 把肩关节推到限位（占位惯量与开环条件下）。

\begin{note}[陷阱：把限位夹持误诊为重力补偿出错]
arm\_dm 调试中曾出现「机械臂在 home 位不悬浮、肩关节往下垂」的现象，第一反应是 $\mathbf{g}(\mathbf{q})$ 算错或 URDF 质量错。逐项核验后发现：URDF 与仿真模型的连杆质量\emph{完全一致}、Pinocchio 的 $\mathbf{g}(\mathbf{q})$ 也正确。真因是\emph{运动学的}：肩关节起始恰在其上限 $0\,\mathrm{rad}$，重力把它\emph{推向限位}、被限位\emph{夹住}不动；给一个 $-6\,\mathrm{N\!\cdot\!m}$ 推离限位，它立刻摆到下限。教训：诊断要\emph{数据驱动}——零力矩自由落体 + 已知力矩权限 + 模型质量比对，三管齐下定位到「限位夹持」，而非盲调增益。这条「先验证通路、再怀疑算法」的方法论贯穿全章故障诊断。
\end{note}

\paragraph{反面：调节器跟不动时变轨迹。}
PD+重力补偿的无源性证明（\cref{thm:ctrl-pd-grav}）有一个\emph{致命的前提}：$\mathbf{q}_d$ 是\emph{常值}（点到点定位）。一旦目标变成时变轨迹 $\mathbf{q}_d(t)$，把 $\boldsymbol{\tau}=\mathbf{g}(\mathbf{q})+\mathbf{K}_p(\mathbf{q}_d(t)-\mathbf{q})-\mathbf{K}_d\dot{\mathbf{q}}$ 代入动力学 \cref{eq:actrl-manip}，闭环里多出两项无人抵消的「干扰」——惯性项 $\mathbf{M}\ddot{\mathbf{q}}_d$ 与科氏项 $\mathbf{C}\dot{\mathbf{q}}_d$。它们随轨迹的加速度/速度起伏，PD 只能\emph{事后}纠偏，于是出现\emph{周期性滞后}：误差波形与轨迹同频、相位落后。arm\_dm 的实测里，这个滞后达 $50$--$200\,\mathrm{mrad}$（占位惯量下）。要消掉它，控制律必须把 $\mathbf{M}\ddot{\mathbf{q}}_d$ 与 $\mathbf{C}\dot{\mathbf{q}}_d$ \emph{前馈}进去——这正是下一节计算力矩的全部动机。

% ════════════════════════════════════════════════════════════════
\section{关节空间：计算力矩跟踪}
\label{sec:actrl-ctc}

\paragraph{从「调节」到「跟踪」。}
承 \cref{ins:actrl-decouple-valid}：当独立关节假设失效（直驱/低减速比/高速、耦合不再被 $1/k_r^2$ 压住）时，计算力矩正是那条「用模型替你解耦」的对策——它把满阵 $\mathbf{M}(\mathbf{q})$ 与 $\mathbf{C}(\mathbf{q},\dot{\mathbf{q}})$ 显式算进控制律、主动抵消耦合，从而\emph{不依赖}齿轮箱去解耦。计算力矩控制（computed torque）的控制律、其线性解耦定理（$\ddot{\mathbf{e}}+\mathbf{K}_d\dot{\mathbf{e}}+\mathbf{K}_p\mathbf{e}=\mathbf{0}$）与完整证明，已在 \cref{sec:ctrl-robot-ctc} 与 \cref{thm:ctrl-ctc} 给出，本章不重证。这里把它专门落到\emph{轨迹跟踪}场景，并补上 P4 没有的两件东西：一台真臂上的高效实现，和一组隔离了被测变量的\emph{公平对比}数据。复用 \cref{eq:ctrl-ctc} 并记跟踪误差 $\mathbf{e}=\mathbf{q}_d-\mathbf{q}$、虚拟加速度输入
\begin{equation}
  \mathbf{a}_{\rm cmd}=\ddot{\mathbf{q}}_d+\mathbf{K}_p\mathbf{e}+\mathbf{K}_d\dot{\mathbf{e}},\qquad
  \boxed{\boldsymbol{\tau}=\mathbf{M}(\mathbf{q})\,\mathbf{a}_{\rm cmd}+\mathbf{C}(\mathbf{q},\dot{\mathbf{q}})\dot{\mathbf{q}}+\mathbf{g}(\mathbf{q})=\mathrm{RNEA}(\mathbf{q},\dot{\mathbf{q}},\mathbf{a}_{\rm cmd}).}
  \label{eq:actrl-ctc}
\end{equation}

\begin{insight}[一次逆动力学算尽前馈]
\cref{eq:actrl-ctc} 的右端等式是工程上的关键：递归牛顿--欧拉算法（RNEA，见 \cref{ch:arm-dyn} 的 $O(n)$ 逆动力学）\emph{本身}就是「给定 $(\mathbf{q},\dot{\mathbf{q}},\ddot{\mathbf{q}})$ 算 $\boldsymbol{\tau}$」。把第三个参数喂 $\mathbf{a}_{\rm cmd}$（而非真实 $\ddot{\mathbf{q}}$），\emph{一次}RNEA 调用就同时算出了 $\mathbf{M}\mathbf{a}_{\rm cmd}+\mathbf{C}\dot{\mathbf{q}}+\mathbf{g}$ 整项——\emph{无需}显式构造 $\mathbf{M}$、$\mathbf{C}$、$\mathbf{g}$ 三个矩阵再相乘。这是计算力矩能在 $\mathrm{kHz}$ 级实时跑的根本（构造完整 $\mathbf{M}$ 需 CRBA、是 $O(n^2)$；而这里只要 $O(n)$）。其代价是\emph{需要精确模型}：$\hat{\mathbf{M}},\hat{\mathbf{C}},\hat{\mathbf{g}}$ 有偏，\cref{thm:ctrl-ctc} 的精确抵消就退化为带残差的近似线性化——这正是它与 PD+重力补偿的分工（\cref{tab:ctrl-ctc-vs-pd}）。
\end{insight}

\paragraph{时序的隐蔽要点。}
跟踪律里 $\ddot{\mathbf{q}}_d,\dot{\mathbf{q}}_d$ 是\emph{前馈}量，必须按轨迹的\emph{物理时间}求值，而非按控制回调的到达时刻。实现上应以测量消息自带的时间戳（仿真时间/硬件时间）去查询期望轨迹，使前馈与\emph{回调率无关}；否则回调抖动会把错误的 $\ddot{\mathbf{q}}_d$ 注入 $\mathbf{a}_{\rm cmd}$，破坏前馈精度。

\paragraph{实测：计算力矩比 PD+重力补偿紧 $26\times$。}
在 arm\_dm 上做了一组对照：同一条时变轨迹，分别用计算力矩 \cref{eq:actrl-ctc} 与 PD+重力补偿 \cref{eq:ctrl-pd-grav} 跟踪。结果\emph{计算力矩 RMS $4.1\,\mathrm{mrad}$ vs PD+重力补偿 RMS $107.4\,\mathrm{mrad}$——紧 $26$ 倍}。PD+重力补偿那 $50$--$200\,\mathrm{mrad}$ 的周期滞后，正是计算力矩前馈掉的 $\mathbf{M}\ddot{\mathbf{q}}_d+\mathbf{C}\dot{\mathbf{q}}_d$ 项（见上节反面分析）。这把「前馈逆动力学有什么用」从定理变成了\emph{可测的数字}。

\begin{note}[如何设计一个「公平」的算法对比：隔离被测变量]
这组 $26\times$ 的对比，最值得学的不是数字本身，而是\emph{怎样让它公平}。被测变量应当只有「跟踪律」一项，其余必须等同、且不能让无关自由度污染基线。arm\_dm 的对比做了三件事：
\begin{enumerate}
  \item \textbf{各自全新初始化}：两种律各从同一 home 起点\emph{独立}启动——因为力矩控制器退出后会\emph{保持}最后力矩，run 间隙臂会漂走，复用进程会让起点不公平。
  \item \textbf{增益异质}：取 $\mathbf{K}_p=\mathrm{diag}([150,150,150,150,100,60])$，腕端故意降低以匹配其小惯量（理由见 \cref{sec:actrl-heterogeneous}），临界阻尼 $\mathbf{K}_d=2\sqrt{\mathbf{K}_p}$。
  \item \textbf{只激励高惯量关节，且基线不被腕拖累}：只让良态的肩肘 $\mathrm{J}2/\mathrm{J}3$ 摆动，其余握住；关键是被握住的关节在\emph{两种}模式下\emph{都}用计算力矩的保持律，度量只在受激关节上算。这样既隔离了「跟踪律」本身，又避免让均匀增益 PD 去摇晃 $2\times10^{-4}$ 惯量的腕（否则 PD+重力补偿基线会因腕失稳而\emph{不公平地}变差）。
\end{enumerate}
这是「对比实验设计」的范例：先想清楚\emph{什么必须等同、什么不能污染基线}，再读数字。
\end{note}

\begin{note}[诚实框定：这个 $26\times$ 的边界]
独立复核提醒：上述计算力矩的精确跟踪，部分依赖了仿真层注入的关节阻尼与「腕解耦」（见 \cref{sec:actrl-heterogeneous}），并非教科书纯净的 $\ddot{\mathbf{e}}+\mathbf{K}_d\dot{\mathbf{e}}+\mathbf{K}_p\mathbf{e}=\mathbf{0}$；且惯量含占位成分。这不削弱「前馈消滞后」的结论，但提醒读者：实测数字要\emph{如实标注其条件}（占位惯量、附加阻尼），这是工程素养，也是 \cref{sec:actrl-heterogeneous} 异质惯量臂专精的引子。
\end{note}

% ════════════════════════════════════════════════════════════════
\section{异质惯量臂的控制率选择 \texorpdfstring{$\bigstar$}{(进阶)}}
\label{sec:actrl-heterogeneous}

\begin{quote}\small\itshape
本节是 P4 完全没有的内容，也是本部区别于一般教科书的独有价值：当一台臂的惯量分布\emph{极不均匀}（重肩肘 + 轻腕）时，「该用哪种控制律、增益该怎么设」不再是品味问题，而由\emph{物理}强制决定。难度 $\bigstar\bigstar$（进阶/选读，但强烈建议——它把前两节的「律」与下一节的「数值病」串成因果链）。
\end{quote}

\paragraph{动机：一台「半良态半病态」的臂。}
理想教科书臂常默认惯量均匀、质量阵良态。真臂往往不是：六轴串联臂的典型结构是\emph{重肩肘 + 轻球腕}——腰肩肘 $\mathrm{J}1/\mathrm{J}2/\mathrm{J}3$ 承载整条臂、等效惯量大且良态；而球腕 $\mathrm{J}4/\mathrm{J}5/\mathrm{J}6$ 只驱动末端小法兰，真实 CAD 惯量可低至 $\sim2\times10^{-4}\,\mathrm{kg\!\cdot\!m^2}$。这\emph{两个数量级}的惯量落差，是本节一切问题的物理根源，也直接决定了 \cref{sec:ad-conditioning} 里质量阵为何\emph{固有病态}。

\subsection{为何低惯量关节\emph{更难}控：自然时间常数 vs 控制频率}
\label{sec:actrl-timescale}

\paragraph{反直觉：轻 $\ne$ 好控。}
直觉会说「关节越轻越灵活、越好控」。恰恰相反。把单关节近似为二阶系统 $I\ddot{q}+b\dot{q}+kq=\tau$，其\emph{自然频率}与\emph{自然时间常数}为
\begin{equation}
  \omega_n=\sqrt{k/I},\qquad \tau_{\rm nat}\approx I/b.
  \label{eq:actrl-omegan}
\end{equation}
惯量 $I$ 出现在\emph{分母}：$I$ 越小，$\omega_n$ 越\emph{高}、$\tau_{\rm nat}$ 越\emph{短}。对腕 $I\approx2\times10^{-4}$、阻尼 $b\approx0.1$，自然时间常数 $\tau_{\rm nat}\approx2\,\mathrm{ms}$，对应特征频率约 $500\,\mathrm{Hz}$。

\begin{theorem}[采样定理对控制率的强制]
\label{thm:actrl-undersample}
一个开环对象的最快动态时间常数为 $\tau_{\rm nat}$、其上闭一个采样周期为 $\Delta t$ 的离散反馈环。若控制频率\emph{不足以分辨}该动态（$\Delta t\gtrsim\tau_{\rm nat}$，即每个时间常数内采样不足若干次），离散闭环会\emph{欠采样发散}：反馈在采样间隔内来不及更新，等效相位滞后超过稳定裕度。\rebuilt{（采样控制稳定性的工程判据，据 \cite{slotine1991applied} 离散化稳定性讨论改述）}
\end{theorem}

\begin{insight}
\cref{thm:actrl-undersample} 解释了腕为何发散。腕的自然频率 $\sim500\,\mathrm{Hz}$ \emph{快于}典型控制率（$\sim100$--$200\,\mathrm{Hz}$）——控制环根本\emph{看不清}腕的动态，于是它在控制层面欠采样、高频发散。一个吊诡的推论：让腕\emph{保持稳定}的，恰恰是仿真物理层（或真机机械层）那个跑在 $500\,\mathrm{Hz}$ 以上的隐式关节阻尼；而这个阻尼是\emph{物理层的}，控制器\emph{无法在自己的层面补偿它}。这就引出了本节的核心论断：异质惯量臂的低惯量关节，在常规控制率下\emph{不能}用大范围闭环跟踪去「硬控」，必须另想办法（解耦）。
\end{insight}

\begin{note}[实测：对 $2\times10^{-4}$ 惯量腕施原始 PD 的灾难]
arm\_dm 给出了最有说服力的量化实证。对那只 $I\approx2\times10^{-4}$ 的腕施加\emph{均匀增益}的原始 PD（$\mathbf{K}_p=150$）：其自然频率 $\omega_n=\sqrt{150/2\times10^{-4}}\approx866\,\mathrm{rad/s}\approx138\,\mathrm{Hz}$ 的闭环带宽要求远超控制率，于是\emph{力矩在钳位前冲到 $888\,\mathrm{N\!\cdot\!m}$}（对一只 $26\,\mathrm{N\!\cdot\!m}$ 额定的腕！），\emph{即便钳位到 $26\,\mathrm{N\!\cdot\!m}$ 后仍产生 $1.3\times10^5\,\mathrm{rad/s^2}$ 的加速度发散}。一个曾试过的「修法」是加黏滞前馈 $\boldsymbol{\tau}\mathrel{+}=0.1\dot{\mathbf{q}}$ 想让模型匹配对象，结果\emph{更糟}（它抵消掉了物理层那个救命的稳定阻尼），RMS 反升——只能撤回。结论：低惯量末端在控制率下「原始 PD 必炸、不可 PD 控」，是设计约束而非可调参数。
\end{note}

\subsection{解药一：增益异质——让 \texorpdfstring{$\mathbf{M}(\mathbf{q})$}{M(q)} 给出正确的力矩权限}
\label{sec:actrl-gainhetero}

\paragraph{均匀增益的错配。}
\cref{thm:ctrl-ctc} 说计算力矩使闭环\emph{解耦}为 $n$ 个独立二阶环 $\ddot{e}_j+k_{d,j}\dot{e}_j+k_{p,j}e_j=0$，但这是\emph{误差}动态。落到\emph{力矩}时，控制律 $\boldsymbol{\tau}=\mathbf{M}(\mathbf{q})\mathbf{a}_{\rm cmd}+\cdots$ 把虚拟加速度 $\mathbf{a}_{\rm cmd}=\mathbf{K}_p\mathbf{e}+\cdots$ 经\emph{质量阵 $\mathbf{M}$ 加权}成力矩。这正是计算力矩比均匀增益 PD 优越的\emph{第二个}理由（第一个是消滞后）：

\begin{insight}[计算力矩天然给每关节「正确的力矩权限」]
均匀增益 PD 对所有关节用同一个 $\mathbf{K}_p$，于是给重肩肘和轻腕\emph{同样大}的力矩响应——对重关节合适，对轻腕则\emph{过激}（同样的位置误差，轻腕只需极小力矩就能纠正，大增益会让它过冲发散，正是上面 $888\,\mathrm{N\!\cdot\!m}$ 的来历）。而计算力矩里的 $\mathbf{M}(\mathbf{q})$ \emph{加权}：重关节的 $M_{jj}$ 大，同样的 $a_{\rm cmd}$ 配大力矩；轻腕的 $M_{jj}$ 小，自动配\emph{小}力矩。换言之，$\mathbf{M}(\mathbf{q})$ 把「需要多大力气推动这个关节」的物理信息\emph{编进}了控制律——这是用逆动力学控制异质惯量臂的根本优势。
\end{insight}

\paragraph{工程对策：异质 $\mathbf{K}_p$。}
即便用计算力矩，外环增益 $\mathbf{K}_p$ 仍应\emph{异质设置}以匹配各关节的有效带宽：arm\_dm 取 $\mathbf{K}_p=\mathrm{diag}([150,150,150,150,100,60])$，腕端 $\mathrm{J}5/\mathrm{J}6$ 故意降到 $100/60$，临界阻尼 $\mathbf{K}_d=2\sqrt{\mathbf{K}_p}$（逐分量）。这把外环的目标闭环频率压到控制率分辨得了的范围。

\subsection{解药二：解耦低惯量腕——切断 \texorpdfstring{$\mathbf{M}$}{M} 耦合腐蚀}
\label{sec:actrl-decouple}

\paragraph{为何增益异质还不够：耦合腐蚀。}
仅降腕增益，在\emph{小范围}激励下够用（\cref{sec:actrl-ctc} 的 $26\times$ 对比即如此）；但在\emph{大范围重定位}时还会暴露第二重病。关键在 $\mathbf{a}_{\rm cmd}=\mathbf{K}_p\mathbf{e}$ 是对\emph{全部 $n$ 个关节}的反馈，而力矩 $\boldsymbol{\tau}=\mathbf{M}(\mathbf{q})\mathbf{a}_{\rm cmd}+\cdots$ 里的 $\mathbf{M}$ 是\emph{满阵}（关节间惯量耦合）。于是某关节 $j$ 的力矩
\begin{equation}
  \tau_j=\sum_{k}M_{jk}(\mathbf{q})\,a_{{\rm cmd},k}+(\mathbf{C}\dot{\mathbf{q}})_j+g_j
  \label{eq:actrl-coupling}
\end{equation}
含\emph{所有}关节的 $a_{{\rm cmd},k}$。若腕（$k=$ 腕）误差 $e_k$ 因欠采样而漂大，它经\emph{交叉惯量项} $M_{jk}$ 注入主关节 $j$ 的力矩，把本该 $3.5\,\mathrm{N\!\cdot\!m}$ 的 $\tau_2$ 推到 $15\,\mathrm{N\!\cdot\!m}$——主关节被腕的误差「腐蚀」、跟不上、卡住。

\begin{theorem}[腕误差经质量阵耦合腐蚀主关节]
\label{thm:actrl-corruption}
设腕关节因控制率不足而误差不可忽略（$a_{{\rm cmd,wrist}}\ne0$）。则由 \cref{eq:actrl-coupling}，主关节力矩 $\tau_j$ 含腐蚀项 $\sum_{k\in{\rm wrist}}M_{jk}(\mathbf{q})\,a_{{\rm cmd},k}$。\textbf{解耦律}：将腕的虚拟加速度\emph{置零}
\begin{equation}
  a_{{\rm cmd},k}=0,\quad k\in{\rm wrist},
  \label{eq:actrl-decouple}
\end{equation}
则计算 $\tau_j$ 时腕加速度按 $0$ 计入 RNEA，腐蚀项消失，主关节力矩还原；腕力矩另以温和保持律（小 $\mathbf{K}_p$ + 物理层高阻尼兜底）施加。
\end{theorem}

\begin{insight}
\cref{thm:actrl-corruption} 的解耦律物理上完全自洽：腕在大范围重定位中\emph{近乎静止}（路径常锁腕），其真实加速度本就 $\approx0$，故把 $a_{{\rm cmd,wrist}}$ 置零\emph{不引入模型误差}，反而剔除了「用一个错的腕加速度去算主关节力矩」的污染。arm\_dm 的对比验证：去耦合后 $\tau_2$ 回到 $\sim3.5\,\mathrm{N\!\cdot\!m}$（与小范围基线一致）、精确跟踪。这把 \cref{sec:actrl-ctc} 的「只动高惯量关节、握住腕」在大范围场景升级成「\emph{显式解耦}腕（零其 $a_{\rm cmd}$ + 阻尼兜底）」。一条可迁移的设计原则：\textbf{异质惯量臂做大范围计算力矩跟踪，必须把低惯量自由度从反馈链里解耦出去}，否则其误差会经满质量阵反噬主关节。
\end{insight}

\begin{algo}[异质惯量臂的计算力矩跟踪（含解耦）]
\label{algo:actrl-hetero-ct}
给定期望轨迹 $(\mathbf{q}_d,\dot{\mathbf{q}}_d,\ddot{\mathbf{q}}_d)(t)$、异质增益 $\mathbf{K}_p,\mathbf{K}_d$、主关节集 $\mathcal{A}$、腕集 $\mathcal{W}$：
\begin{enumerate}
  \item 读状态 $(\mathbf{q},\dot{\mathbf{q}})$，按\emph{消息时间戳}查 $(\mathbf{q}_d,\dot{\mathbf{q}}_d,\ddot{\mathbf{q}}_d)$；
  \item 误差 $\mathbf{e}=\mathbf{q}_d-\mathbf{q},\ \dot{\mathbf{e}}=\dot{\mathbf{q}}_d-\dot{\mathbf{q}}$；
  \item 虚拟加速度：$a_{{\rm cmd},j}=\ddot{q}_{d,j}+k_{p,j}e_j+k_{d,j}\dot{e}_j$（$j\in\mathcal{A}$）；$\;a_{{\rm cmd},k}=0$（$k\in\mathcal{W}$，解耦）；
  \item 力矩：$\boldsymbol{\tau}=\mathrm{RNEA}(\mathbf{q},\dot{\mathbf{q}},\mathbf{a}_{\rm cmd})$；腕力矩另叠温和保持 $\tau_k\mathrel{+}=k_p^{\rm hold}(q_{d,k}-q_k)$；
  \item 钳位 $\boldsymbol{\tau}\leftarrow\mathrm{clip}(\boldsymbol{\tau},\pm\tau_{\max})$，下发。
\end{enumerate}
\end{algo}

% ════════════════════════════════════════════════════════════════
\section{操作空间控制 OSC（完整推导）}
\label{sec:actrl-osc}

\paragraph{动机：在末端坐标里直接施力。}
\cref{sec:ctrl-robot-osc} 给了 Khatib 操作空间公式的「一张图景」：任务空间动力学 \cref{eq:ctrl-osc} 与对偶映射 $\boldsymbol{\tau}=\mathbf{J}^\top\mathbf{F}$。本节把它在串联臂上\emph{完整推导}——从关节动力学一步步投影到任务空间、导出操作空间惯量、再导出\emph{动力学一致}的零空间投影——并指出它对质量阵条件数的高度敏感（这正是 \cref{sec:actrl-heterogeneous} 病态故事的任务空间版）。P4 的 \cref{ch:operational-space} 把这套理论展开到了浮动基与稳定性证明，本节与之互链而不重复，只给固定基串联臂的自包含版本。理论据 Khatib\cite{khatib1987operational}，公式经网络核对\rebuilt{}。

\subsection{从关节动力学到任务空间动力学}
\label{sec:actrl-osc-deriv}

设末端任务坐标 $\mathbf{x}\in\mathbb{R}^m$、任务雅可比 $\mathbf{J}(\mathbf{q})\in\mathbb{R}^{m\times n}$（$\dot{\mathbf{x}}=\mathbf{J}\dot{\mathbf{q}}$，见 \cref{ch:arm-kin}）。对其求导得任务加速度
\begin{equation}
  \ddot{\mathbf{x}}=\mathbf{J}\ddot{\mathbf{q}}+\dot{\mathbf{J}}\dot{\mathbf{q}}.
  \label{eq:actrl-osc-xddot}
\end{equation}
从关节动力学 \cref{eq:actrl-manip}（无外力）解出 $\ddot{\mathbf{q}}=\mathbf{M}^{-1}(\boldsymbol{\tau}-\mathbf{C}\dot{\mathbf{q}}-\mathbf{g})$，代入 \cref{eq:actrl-osc-xddot}：
\begin{equation}
  \ddot{\mathbf{x}}=\mathbf{J}\mathbf{M}^{-1}(\boldsymbol{\tau}-\mathbf{C}\dot{\mathbf{q}}-\mathbf{g})+\dot{\mathbf{J}}\dot{\mathbf{q}}.
  \label{eq:actrl-osc-step}
\end{equation}
设关节力矩由对偶映射 $\boldsymbol{\tau}=\mathbf{J}^\top\mathbf{F}$ 产生（$\mathbf{F}$ 为末端广义力，对偶性由虚功原理 $\boldsymbol{\tau}^\top\delta\mathbf{q}=\mathbf{F}^\top\delta\mathbf{x},\ \delta\mathbf{x}=\mathbf{J}\delta\mathbf{q}$ 保证，见 \cref{ch:arm-kin} 静力学对偶）。代入并左乘 $(\mathbf{J}\mathbf{M}^{-1}\mathbf{J}^\top)^{-1}$，可解出 $\mathbf{F}$ 与 $\ddot{\mathbf{x}}$ 的关系。定义\textbf{操作空间惯量矩阵}
\begin{equation}
  \boxed{\boldsymbol{\Lambda}(\mathbf{q})=(\mathbf{J}\mathbf{M}^{-1}\mathbf{J}^\top)^{-1}\in\mathbb{R}^{m\times m},}
  \label{eq:actrl-Lambda}
\end{equation}
整理 \cref{eq:actrl-osc-step} 即得\textbf{任务空间（操作空间）动力学}
\begin{equation}
  \boxed{\boldsymbol{\Lambda}(\mathbf{q})\ddot{\mathbf{x}}+\boldsymbol{\mu}(\mathbf{q},\dot{\mathbf{q}})+\mathbf{p}(\mathbf{q})=\mathbf{F},}
  \label{eq:actrl-osc-dyn}
\end{equation}
其中任务空间科氏/离心 $\boldsymbol{\mu}=\bar{\mathbf{J}}^\top\mathbf{C}\dot{\mathbf{q}}-\boldsymbol{\Lambda}\dot{\mathbf{J}}\dot{\mathbf{q}}$、任务空间重力 $\mathbf{p}=\bar{\mathbf{J}}^\top\mathbf{g}$，$\bar{\mathbf{J}}$ 即下面的动力学一致逆。\cref{eq:actrl-osc-dyn} 与关节空间动力学 \cref{eq:actrl-manip} \emph{形式同构}：$\boldsymbol{\Lambda}$ 是末端「感受到的」等效惯量。

\subsection{动力学一致广义逆与零空间投影}
\label{sec:actrl-osc-null}

\paragraph{冗余的自由度去哪了。}
当 $n>m$（冗余臂，或如 arm\_dm 把 $6$ 维位姿任务降维到 $3$ 维位置任务而留出冗余），$\mathbf{J}^\top\in\mathbb{R}^{n\times m}$ 不满秩列空间之外还有 $n-m$ 维\emph{零空间}：存在关节力矩 $\boldsymbol{\tau}_0$ 满足 $\mathbf{J}\mathbf{M}^{-1}\boldsymbol{\tau}_0$ 不一定为零——能否找一个投影，让任意 $\boldsymbol{\tau}_0$ 投进零空间后\emph{绝不}扰动末端加速度？这正是「动力学一致」的要义。

\begin{definition}[动力学一致广义逆]
\label{def:actrl-jbar}
$\mathbf{J}^\top$ 的\textbf{动力学一致广义逆}定义为
\begin{equation}
  \boxed{\bar{\mathbf{J}}=\mathbf{M}^{-1}\mathbf{J}^\top\boldsymbol{\Lambda}=\mathbf{M}^{-1}\mathbf{J}^\top(\mathbf{J}\mathbf{M}^{-1}\mathbf{J}^\top)^{-1},\qquad \bar{\mathbf{J}}^\top=\boldsymbol{\Lambda}\mathbf{J}\mathbf{M}^{-1}.}
  \label{eq:actrl-jbar}
\end{equation}
它是\emph{以 $\mathbf{M}$ 为度量}的加权伪逆（而非 Moore--Penrose 的 $\mathbf{J}^+=\mathbf{J}^\top(\mathbf{J}\mathbf{J}^\top)^{-1}$，后者以单位阵为度量）。
\end{definition}

\begin{theorem}[零空间力矩对末端零加速度]
\label{thm:actrl-null}
\textbf{动力学一致零空间投影}
\begin{equation}
  \boxed{\mathbf{N}=\mathbf{I}-\mathbf{J}^\top\bar{\mathbf{J}}^\top=\mathbf{I}-\mathbf{J}^\top(\mathbf{J}\mathbf{M}^{-1}\mathbf{J}^\top)^{-1}\mathbf{J}\mathbf{M}^{-1}\in\mathbb{R}^{n\times n}}
  \label{eq:actrl-N}
\end{equation}
使任意零空间力矩 $\mathbf{N}\boldsymbol{\tau}_0$ 对末端产生\emph{零}加速度。完整控制律
\begin{equation}
  \boldsymbol{\tau}=\mathbf{J}^\top\mathbf{F}+\mathbf{N}\boldsymbol{\tau}_0,\qquad
  \mathbf{F}=\hat{\boldsymbol{\Lambda}}\mathbf{F}^*+\hat{\boldsymbol{\mu}}+\hat{\mathbf{p}},\quad
  \mathbf{F}^*=\ddot{\mathbf{x}}_d+\mathbf{K}_d(\dot{\mathbf{x}}_d-\dot{\mathbf{x}})+\mathbf{K}_p(\mathbf{x}_d-\mathbf{x}),
  \label{eq:actrl-osc-law}
\end{equation}
其中 $\mathbf{F}^*$ 是力级 PD，给出解耦行为 $\ddot{\mathbf{x}}=\mathbf{F}^*$；$\boldsymbol{\tau}_0$ 为次任务（如关节中心化、避奇异）力矩。\footnote{作用在关节力矩上的投影取 $\mathbf{N}\boldsymbol{\tau}_0$（左乘 $\mathbf{N}$、不转置），与 Khatib\cite{khatib1987operational} 的标准式 $\boldsymbol{\tau}=\mathbf{J}^\top\mathbf{F}+(\mathbf{I}-\mathbf{J}^\top\bar{\mathbf{J}}^\top)\boldsymbol{\tau}_0$ 一致；由于 $\mathbf{N}$ 非对称，$\mathbf{N}^\top$ 一侧不满足 $\mathbf{J}\mathbf{M}^{-1}(\cdot)=\mathbf{0}$，见下证。}
\end{theorem}
\begin{proof}\rebuilt{（据 \cite{khatib1987operational}，网络核对）}
末端加速度由 \cref{eq:actrl-osc-xddot} 与 $\ddot{\mathbf{q}}=\mathbf{M}^{-1}(\boldsymbol{\tau}-\mathbf{C}\dot{\mathbf{q}}-\mathbf{g})$ 给出，其对 $\boldsymbol{\tau}$ 的依赖为 $\mathbf{J}\mathbf{M}^{-1}\boldsymbol{\tau}$。故只需验证零空间分量 $\mathbf{N}\boldsymbol{\tau}_0$ 经 $\mathbf{J}\mathbf{M}^{-1}$ 映射后为零。由 \cref{eq:actrl-N} 的 $\mathbf{N}=\mathbf{I}-\mathbf{J}^\top\bar{\mathbf{J}}^\top=\mathbf{I}-\mathbf{J}^\top\boldsymbol{\Lambda}\mathbf{J}\mathbf{M}^{-1}$（代入 $\bar{\mathbf{J}}^\top=\boldsymbol{\Lambda}\mathbf{J}\mathbf{M}^{-1}$），左乘 $\mathbf{J}\mathbf{M}^{-1}$ 并用 $\mathbf{J}\mathbf{M}^{-1}\mathbf{J}^\top=\boldsymbol{\Lambda}^{-1}$（由 \cref{eq:actrl-Lambda}）：
\[
  \mathbf{J}\mathbf{M}^{-1}\mathbf{N}
  =\mathbf{J}\mathbf{M}^{-1}-\underbrace{\mathbf{J}\mathbf{M}^{-1}\mathbf{J}^\top}_{=\,\boldsymbol{\Lambda}^{-1}}\boldsymbol{\Lambda}\,\mathbf{J}\mathbf{M}^{-1}
  =\mathbf{J}\mathbf{M}^{-1}-\boldsymbol{\Lambda}^{-1}\boldsymbol{\Lambda}\,\mathbf{J}\mathbf{M}^{-1}=\mathbf{0}.
\]
中间只有\emph{一个} $\mathbf{M}^{-1}$ 被 $\mathbf{J}\cdots\mathbf{J}^\top$ 夹成 $\boldsymbol{\Lambda}^{-1}$ 而精确消去，这正是把投影矩阵作用在\emph{力矩}上（左乘 $\mathbf{N}$，而非 $\mathbf{N}^\top$）的关键：$\mathbf{N}$ 一般\emph{非对称}（$\mathbf{N}^\top=\mathbf{I}-\mathbf{M}^{-1}\mathbf{J}^\top\boldsymbol{\Lambda}\mathbf{J}\ne\mathbf{N}$），唯有 $\mathbf{N}$ 这一侧使任意 $\boldsymbol{\tau}_0$ 对末端\emph{零}加速度。$\mathbf{N}$ 还是\emph{幂等}投影（$\mathbf{N}^2=\mathbf{N}$，本节练习），其值域恰为零空间。$\bar{\mathbf{J}}$ 的\emph{唯一性}由「以 $\mathbf{M}$ 为度量的最小动能解」刻画：它是使 $\tfrac12\boldsymbol{\tau}_0^\top\mathbf{M}^{-1}\boldsymbol{\tau}_0$ 最小的逆。
\end{proof}

\begin{insight}[为何不能用 Moore--Penrose 伪逆]
许多工业实现（如某些笛卡尔阻抗示例）图省事，把零空间投影写成纯运动学的 $(\mathbf{I}-\mathbf{J}^\top\mathbf{J}^{+\top})$（用 $\mathbf{J}^+$ 而非 $\bar{\mathbf{J}}$）。这\emph{不是}动力学一致的：它以单位阵为度量，忽略了 $\mathbf{M}$。后果——零空间里的关节中心化/避奇异力矩会「\emph{泄漏}」出一个末端寄生加速度，破坏任务的纯净（末端被次任务\emph{推动}）。\cref{thm:actrl-null} 的 $\bar{\mathbf{J}}$ 是\emph{唯一}让这种泄漏归零的逆，代价是它含 $\mathbf{M}^{-1}$——这把 OSC 的成败\emph{绑死}在了质量阵的数值条件上（下小节）。这条「两种伪逆务必严格区分」的工程缘由，也是 \cref{ch:force-wbc} 反复强调的。
\end{insight}

\subsection{OSC 对质量阵条件数的敏感性 \texorpdfstring{$\bigstar$}{(进阶)}}
\label{sec:actrl-osc-illcond}

\paragraph{动力学一致逆的阿喀琉斯之踵。}
\cref{eq:actrl-jbar,eq:actrl-N} 里那个 $\mathbf{M}^{-1}$ 不是免费的。当质量阵\emph{病态}（条件数 $\mathrm{cond}(\mathbf{M})$ 大，正是 \cref{sec:ad-conditioning} 异质惯量臂的固有处境）时，$\mathbf{M}^{-1}$ 数值上被极小特征值放大，$(\mathbf{J}\mathbf{M}^{-1}\mathbf{J}^\top)^{-1}$ 在近奇异处更雪上加霜。于是\emph{动力学一致投影本应有的「特征值非 $0$ 即 $1$」的纯投影性质被数值误差破坏}——投影矩阵 $\mathbf{N}$ 不再幂等，零空间力矩的「泄漏」复活，末端漂移。

\begin{note}[实测：异质惯量臂的 OSC 零空间收不住末端（病态机理的穷尽验证）]
arm\_dm 是六轴臂，把 $6$ 维位姿任务\emph{降维}到 $3$ 维末端位置，留 $6-3=3$ 维零空间做自运动（绕固定工具重构姿态）。期望「末端锁死 + 零空间自运动」，结果\emph{自运动出现了}（肩肘重构 $1.1$--$1.4\,\mathrm{rad}$）但\emph{末端没收住}——漂 $60$--$110\,\mathrm{mm}$ 而非 $<10\,\mathrm{mm}$。穷尽四次定位：
\begin{enumerate}
  \item 用运动学投影 $\mathbf{I}-\mathbf{J}^+\mathbf{J}$（非动力学一致）：末端泄漏 $85\,\mathrm{mm}$；
  \item 用精确动力学投影 \cref{eq:actrl-N}：$\|\boldsymbol{\tau}_0\|$ \emph{爆炸}到 $\sim1.8\times10^4$（腕 $2\times10^{-4}$ 惯量使 $\mathbf{M}^{-1}$ 巨大、奇异处 $(\mathbf{J}\mathbf{M}^{-1}\mathbf{J}^\top)^{-1}$ 又爆）；
  \item 正则化 $\mathbf{M}+\lambda\mathbf{I}$ 或阻尼逆：稳了，但末端仍漂 $70\,\mathrm{mm}$——因正则化\emph{破坏了投影性}（特征值不再 $\{0,1\}$）；
  \item 锁腕：$107\,\mathrm{mm}$，无改善。
\end{enumerate}
真根因经第一性原理坐实：腕惯量 $2\times10^{-4}$ 使 $\mathrm{cond}(\mathbf{M})\approx527$--$1625$（视构型）。即便换上更合理的夹爪惯量把条件数砍到 $300+$，末端\emph{仍漂 $62\,\mathrm{mm}$ 且自运动塌缩}（肩肘只动 $0.07/0.24\,\mathrm{rad}$）——「漏但动」换成了「稳但几乎不动且仍偏」，两者\emph{都不干净}。
\end{note}

\begin{insight}[第一性原理：用条件数判断算法可行性，而非在力矩层硬刚]
这个案例的元教训值千金：\textbf{异质惯量臂的冗余解析正解，是运动学/速度级（无 $\mathbf{M}^{-1}$），而非力矩级 OSC}。同一台臂、同一个零空间任务，换成不含质量阵逆的运动学 resolved-rate（见 \cref{ch:arm-kin} 与 \cref{ch:arm-prac}），末端立刻锁到 $0.00\,\mathrm{mm}$。OSC 不是「更高级所以更好」——它对 $\mathrm{cond}(\mathbf{M})$ \emph{极度敏感}，只在\emph{良态}臂上可行。正确的工程姿态是：先用第一性原理（条件数、特征值结构）\emph{判断}某算法在本臂是否可行，而不是在力矩层反复调正则化「硬刚」一个数值上注定失败的方案。这条判据与 \cref{sec:actrl-heterogeneous} 的「自然时间常数 vs 控制率」并列，构成异质惯量臂控制率选择的两根支柱。
\end{insight}

% ════════════════════════════════════════════════════════════════
\section{阻抗与导纳控制}
\label{sec:actrl-impedance}

\paragraph{动机：接触时纯位置控制会「打架」。}
机械臂插轴、打磨、开门、与人协作时要\emph{接触}环境。纯刚性位置控制（计算力矩/OSC PD）下，微小位置误差被高刚度放大成巨大接触力，卡死或损坏。\cref{sec:ctrl-robot-impedance} 已据 Hogan\cite{hogan1985impedance} 给出核心洞见与 \cref{def:ctrl-impedance}（阻抗 = 运动$\to$力的算子、导纳 = 其逆，且互连双方不能都是导纳，故机械手呈阻抗、环境呈导纳）。本节把它在\emph{力矩控制串联臂}上深化：整定、临界阻尼设计、「真阻抗无需 F/T」范式，并以实测验证线性弹簧律。P4 的 \cref{ch:operational-space} 给了无源性与能量端口的完整证明，本节互链不重复。

\subsection{阻抗 vs 导纳：两种实现的硬件分野}
\label{sec:actrl-imp-vs-adm}

\begin{table}[H]\centering\small
\caption{阻抗、导纳、力位混合三范式}
\label{tab:actrl-imp-adm}
\begin{tabular}{p{0.14\textwidth} p{0.24\textwidth} p{0.24\textwidth} p{0.24\textwidth}}
\toprule
 & 阻抗 impedance & 导纳 admittance & 力位混合 hybrid \\
\midrule
测$\to$出 & 测位置/速度 $\to$ 出\emph{力(力矩)} & 测\emph{力(F/T)} $\to$ 出位置参考 & 选择矩阵\emph{分方向} \\
硬件 & \textbf{力矩臂}（本章），\emph{无需} F/T & 位置臂 + F/T 传感器 & 位置/力轴混合 \\
律 & $\boldsymbol{\tau}=\mathbf{J}^\top(-\mathbf{K}_c\Delta\mathbf{x}-\mathbf{D}_c\dot{\mathbf{x}})+\mathbf{g}$ & $\ddot{\mathbf{x}}=\mathbf{M}_d^{-1}(\mathcal{F}_{\rm ext}-\cdots)$ & Raibert--Craig 选择矩阵 $\mathbf{S}$ \\
带宽/特点 & 高带宽、接触友好、简单、稳 & 高刚度跟踪好、但接触易不稳、带宽受内环限 & 某轴控位某轴控力（打磨/装配）\\
\bottomrule
\end{tabular}
\end{table}

\begin{insight}[「力矩臂 $\to$ 真阻抗」无需力传感器]
为什么一台\emph{力矩控制}的臂能呈现「弹簧」行为而\emph{不装}任何 F/T 传感器？因为阻抗控制是「\emph{测运动、出力}」：它\emph{不读}外力，而是\emph{主动施加}一个「位置偏差 $\to$ 力矩」的弹簧律 $\boldsymbol{\tau}=\mathbf{J}^\top(-\mathbf{K}_c\Delta\mathbf{x})+\cdots$。当环境推动末端、产生位置偏差 $\Delta\mathbf{x}$ 时，这个律\emph{自动}产生反向恢复力矩——末端\emph{表现得}像弹簧，外力越大、退让越多。它需要的全部是\emph{运动学}（雅可比 $\mathbf{J}$ 把笛卡尔力映到关节力矩）与\emph{动力学}（重力 $\mathbf{g}$ 补偿，使弹簧不被自重拖偏），二者皆可由模型算出，\emph{不需}测力。这是「纯力矩驱动臂 $\to$ 真阻抗」范式的精髓，也是它优于「位置臂 + F/T 导纳」的根本（无传感器延迟、高带宽、接触天然稳）。
\end{insight}

\subsection{笛卡尔阻抗律与临界阻尼整定}
\label{sec:actrl-cartimp}

\paragraph{控制律。}
设末端位置误差 $\Delta\mathbf{x}=\mathbf{x}-\mathbf{x}_d$、笛卡尔刚度 $\mathbf{K}_c\succ0$、阻尼 $\mathbf{D}_c\succ0$，\textbf{笛卡尔阻抗控制律}
\begin{equation}
  \boxed{\boldsymbol{\tau}=\mathbf{J}^\top\big(-\mathbf{K}_c\Delta\mathbf{x}-\mathbf{D}_c\dot{\mathbf{x}}\big)+\mathbf{g}(\mathbf{q}),}
  \label{eq:actrl-cartimp}
\end{equation}
即「笛卡尔虚拟弹簧--阻尼」经 $\mathbf{J}^\top$ 落到关节力矩，叠加重力补偿。它是 \cref{eq:ctrl-impedance-law} 免力传感特例在「弹簧--阻尼读法」下的本来面目（取 $\mathbf{M}_d=\boldsymbol{\Lambda}$、不整形惯性）。

\paragraph{整定：临界阻尼。}
要末端\emph{无超调}地回到平衡，阻尼应取\emph{临界}。对一维等效二阶系统 $m_{\rm eff}\ddot{\Delta x}+D\dot{\Delta x}+K\Delta x=0$，临界阻尼条件 $D=2\sqrt{K\,m_{\rm eff}}$。若把末端等效惯量近似为单位（或在 arm\_dm 的工程简化里直接取标量刚度），常用经验整定
\begin{equation}
  \mathbf{D}_c=2\sqrt{\mathbf{K}_c}\quad(\text{逐分量，} m_{\rm eff}\!\approx\!1).
  \label{eq:actrl-critdamp}
\end{equation}

\begin{note}[陷阱：每个笛卡尔方向都临界阻尼时，$\mathbf{D}_c=2\sqrt{\mathbf{K}_c}$ 还不够]
\cref{eq:actrl-critdamp} 隐含了「末端等效惯量是单位阵」的近似。真相是：末端在不同笛卡尔方向「感受到的」惯量是\emph{操作空间惯量} $\boldsymbol{\Lambda}(\mathbf{q})$（\cref{eq:actrl-Lambda}），它\emph{随构型变化、且各向异性}。要在\emph{每个}方向都精确临界阻尼，阻尼矩阵不能是常值，必须\emph{随 $\boldsymbol{\Lambda}$ 整形}（典型做法是「双对角化」$\mathbf{K}_c$ 与 $\boldsymbol{\Lambda}$ 后逐特征方向配 $2\sqrt{\cdot}$）\rebuilt{（据笛卡尔阻抗阻尼设计文献\cite{albu2003cartesian}）}。\cref{eq:actrl-critdamp} 是工程上的够用近似（尤其 arm\_dm 这类标量刚度场景），但读者须知它的边界：精确临界阻尼是\emph{构型相关}的。
\end{note}

\paragraph{平衡处的退化：纯重力补偿。}
\cref{eq:actrl-cartimp} 在末端\emph{停于平衡}（$\Delta\mathbf{x}=\mathbf{0},\dot{\mathbf{x}}=\mathbf{0}$）时退化为 $\boldsymbol{\tau}=\mathbf{g}(\mathbf{q})$——弹簧项为零、只剩重力补偿，臂「失重悬浮」于目标位姿。arm\_dm 实测：HOLD 时 $\|\Delta\mathbf{x}\|=0.0\,\mathrm{mm}$（弹簧项=0、$\boldsymbol{\tau}=$ 纯 $\mathbf{g}$）；追踪 $\pm8\,\mathrm{cm}$ 振荡目标（WIGGLE）时 $\|\Delta\mathbf{x}\|\approx30$--$36\,\mathrm{mm}$ 滞后（有限刚度的柔顺，跟得「软」）。

\subsection{扰动柔顺：线性弹簧律的实证}
\label{sec:actrl-disturb}

\paragraph{把刚度「称」出来。}
阻抗控制的刚度 $\mathbf{K}_c$ 是个\emph{可验证}的物理量：给末端施一个已知外力 $F$，稳态退让 $\Delta x$ 应满足 $\Delta x=F/K$（线性弹簧律）。

\begin{note}[实测：$50\,\mathrm{N}\to0.2\,\mathrm{m}=F/K$ 精确吻合]
arm\_dm 在仿真里给末端连杆施世界系外力 $50\,\mathrm{N}$（取 $\mathbf{K}_c=250\,\mathrm{N/m}$），末端退让 $\|\Delta\mathbf{x}\|=199.9\,\mathrm{mm}\approx F/K=50/250=0.20\,\mathrm{m}$——\emph{教科书级吻合}。撤力后弹回 $1.2\to0.1\,\mathrm{mm}$（线性弹簧 + 临界阻尼，无振荡）。全程\emph{纯力矩、无 F/T}。这把「阻抗控制的刚度真的是一根弹簧」从公式变成了可量的实证；同时也是诚实框定的范例：该力为\emph{开环}施加，验证的是控制律呈现的弹簧律本身。
\end{note}

% ════════════════════════════════════════════════════════════════
\section{力控与力位混合控制}
\label{sec:actrl-hybrid}

\paragraph{动机：某些任务要「分方向」控制。}
打磨要「法向恒力顶住工件、切向自由扫过」；装配要「插入方向控位、侧向控力以顺从对中误差」。阻抗在\emph{所有}方向给同一种「弹簧」行为，而这类任务需要在\emph{不同方向}分别控\emph{位}或控\emph{力}——这是 Raibert--Craig（1981）力位混合控制\cite{raibert1981hybrid} 的用武之地。理论网络重建据\cite{raibert1981hybrid}，公式经核对\rebuilt{}。P4 的 \cref{ch:force-wbc} 把力控谱系（力位混合/主动刚度/并行/阻抗的统一）展开得更全，本节聚焦选择矩阵的核心机理与力矩臂上的工程要点。

\subsection{约束坐标系、自然约束与人工约束}
\label{sec:actrl-constraints}

\paragraph{接触把位置与力「绑」在一起。}
关键观察：在接触方向上，\emph{位置与力不能同时独立指定}——环境几何把二者耦合。沿一面刚性墙，法向\emph{位置}由墙\emph{强制}（你不能穿墙），但法向\emph{力}可自由施加；切向\emph{力}由摩擦\emph{受限}（理想光滑时为零），但切向\emph{位置}可自由移动。这种「环境几何强加的、与任务无关的约束」称\textbf{自然约束}；与之\emph{正交}的、由任务指定的目标称\textbf{人工约束}。

\begin{definition}[约束坐标系与自然/人工约束]
\label{def:actrl-constraint-frame}
在接触处选一\textbf{约束坐标系}（compliance frame），使各坐标轴对齐接触几何（如法向/切向）。在每个轴上：\textbf{自然约束}由环境给定（刚性墙：法向位置固定、切向力为零）；\textbf{人工约束}由任务给定（法向施 $15\,\mathrm{N}$、切向移到某位置）。二者在每个轴上\emph{互补}：自然约束限位置的轴，人工约束指定力；反之亦然。一个含 $N$ 个自然约束与 $N$ 个正交人工约束的坐标系，恰好把 $2N$ 个量（$N$ 位置 + $N$ 力）的指定权\emph{无冲突}地分配完。
\end{definition}

\begin{insight}[为何必须在「精选的」约束坐标系里劈分]
不能随便沿世界系 $x,y,z$ 劈分力/位，因为接触几何（墙面法向、孔轴方向）\emph{一般不对齐}世界轴。只有把约束坐标系\emph{对齐接触几何}，自然约束才在各\emph{单}轴上干净表达（法向轴=位置受限、力自由；切向轴=力受限、位置自由），选择矩阵才能简单地「按轴」取力或取位。选错坐标系，力与位会在同一轴上\emph{纠缠}，混合控制立即失效。这就是 \cref{def:actrl-constraint-frame} 强调「使坐标轴对齐接触几何」的工程缘由。
\end{insight}

\subsection{选择矩阵与混合控制律}
\label{sec:actrl-selection}

\paragraph{把空间劈成两半。}
在约束坐标系里，定义对角\textbf{选择矩阵} $\mathbf{S}=\mathrm{diag}(s_1,\dots,s_m)$，$s_i\in\{0,1\}$：$s_i=1$ 表第 $i$ 轴\emph{控力}，$s_i=0$ 表\emph{控位}。则 $\mathbf{S}$ 取出力子空间、$\mathbf{I}-\mathbf{S}$ 取出位子空间，二者\emph{正交互补}。\textbf{混合控制律}
\begin{equation}
  \boxed{\mathbf{F}=\underbrace{\mathbf{S}\,\mathbf{F}_d}_{\text{力控轴}}+\underbrace{(\mathbf{I}-\mathbf{S})\big(-\mathbf{K}_p\Delta\mathbf{x}-\mathbf{D}\dot{\mathbf{x}}\big)}_{\text{位控轴}},\qquad \boldsymbol{\tau}=\mathbf{J}^\top\mathbf{F}+\mathbf{g}(\mathbf{q}),}
  \label{eq:actrl-hybrid}
\end{equation}
$\mathbf{F}_d$ 为期望力（如法向 $15\,\mathrm{N}$），位控轴用任务空间 PD。力矩仍由对偶映射 $\boldsymbol{\tau}=\mathbf{J}^\top\mathbf{F}$ 落地（叠重力）。

\begin{note}[实测与陷阱：力矩臂上要用全雅可比，不要关节切分]
arm\_dm 做「法向 $x$ 控力（恒 $15\,\mathrm{N}$ 顶墙）+ 切向 $y,z$ 控位」：$\mathbf{S}=\mathrm{diag}(1,0,0)$，$\mathbf{F}=[F_d,\,-K_p(y-y_d)-D\dot{y},\,-K_p(z-z_0)-D\dot{z}]$，$\boldsymbol{\tau}=\mathbf{J}^\top_{(3\times6)}\mathbf{F}+\mathbf{g}$。它经历了\emph{五次}诊断才跑通，三条教训极可迁移：
\begin{enumerate}
  \item \textbf{用全（$6$ 列）雅可比做笛卡尔控制，而非把关节切成「力关节/位关节」}——把冗余的腕安全地放进零空间，避免折叠构型下子雅可比奇异；
  \item \textbf{用关节空间计算力矩去「展开」折叠臂}（笛卡尔直线趋近会撞奇异）；
  \item \textbf{扫掠方向选「无 reach 限制」的自由度}（横向/基座 yaw，等半径不受臂展限），而非被工作空间边界钉死的径向/竖向。
\end{enumerate}
结果：法向 $x$ 钉墙 mean $0.144$、std $3.1\,\mathrm{mm}$；切向 $y$ 跟 $\pm0.05$ 横扫。$\pm22\,\mathrm{mm}$ 的切向滞后\emph{是物理正确的接触摩擦}（$\mu\!\cdot\!15\,\mathrm{N}$ vs $K_p\!\cdot\!0.08$）——正是打磨之所以需要力控之因。
\end{note}

\begin{note}[诚实框定：「钉墙」其实没碰墙]
独立复核指出：上述「钉墙」的受控末端停在离墙 $5.6\,\mathrm{cm}$ 处、\emph{并未真正触墙}；那 $15\,\mathrm{N}$ 是\emph{开环}施加、无任何力反馈，故 std $3\,\mathrm{mm}$ 是开环抖动而非接触摩擦。这条提醒并不否定选择矩阵机理，而是示范工程素养：实验结论要写「占位惯量下、开环力」，删去「钉墙」这类绝对语。诚实记录局限本身就是一种能力。
\end{note}

% ════════════════════════════════════════════════════════════════
\section{故障诊断：力矩跟踪的四种签名}
\label{sec:actrl-diagnostics}

\paragraph{动机：从波形反推根因。}
力矩控制比位置控制更易「炸」，且现象五花八门。一套可复用的诊断法是：把关节时间序列波形按\emph{典型签名}归类，每种签名\emph{大概率}指向某类根因——这把「调参玄学」变成「查表诊断」（这是 v5.0 R15 故障排查在本章的落地）。

\begin{table}[H]\centering\small
\caption{力矩跟踪故障的四种签名 $\to$ 根因映射}
\label{tab:actrl-signatures}
\begin{tabularx}{\linewidth}{p{0.20\textwidth} L L}
\toprule
波形签名 & 长相 & 多半的根因 \\
\midrule
\textbf{stuck-at-home} & 卡在初始位，力矩不饱和却不动 & 限位夹持（\cref{sec:actrl-pdg} 的肩关节顶上限）；或外部几何挡路（碰撞凸包阻挡，用\emph{无障碍基线}隔离验证）\\
\textbf{wrong-sign-runaway} & 朝\emph{一个}方向越跑越快 & 力矩符号反（雅可比/重力符号错）；或 $\mathbf{g}(\mathbf{q})$ 符号/通路错 \\
\textbf{unstable} & 高频抖动直至发散 & 低惯量关节欠采样（\cref{thm:actrl-undersample} 的腕，$\omega_n\gg$ 控制率）；或增益过高、阻尼不足；或 OSC 病态泄漏 \\
\textbf{lagging} & 跟得上但总慢半拍、误差与轨迹同频 & 前馈缺失（PD+重力补偿跟时变轨迹，缺 $\mathbf{M}\ddot{\mathbf{q}}_d+\mathbf{C}\dot{\mathbf{q}}_d$，见 \cref{sec:actrl-pdg} 反面）\\
\bottomrule
\end{tabularx}
\end{table}

\begin{insight}[诊断方法论：先验证数据通道，再怀疑算法]
四签名背后是一条贯穿本章的方法论：\emph{数据驱动诊断，而非盲调}。每次故障先问「是数据/通道问题，还是算法问题」——arm\_dm 多次「看似算法坏」实为通道错（如关节状态\emph{名序错乱}把腕/肩张冠李戴，是 red herring）。配套手段：\textbf{隔离变量法}（跑无障碍基线锁定「挡路」根因）、\textbf{警惕 red herring}（先按\emph{名}而非\emph{下标}读关节状态）、\textbf{诚实框定}（实测标注「占位惯量/开环力」边界）。这套方法论的完整工程展开见 \cref{ch:arm-prac}。
\end{insight}

% ════════════════════════════════════════════════════════════════
\section{常见误解}
\label{sec:actrl-misconceptions}

\begin{enumerate}
  \item \textbf{「机械臂控制就是再调一遍 PID。」}——错。机械臂动力学强非线性、强耦合（\cref{eq:actrl-manip}），且接触任务要控\emph{力/阻抗}而非位置。本章的计算力矩（反馈线性化）、OSC、阻抗、力位混合，都不是单环 PID 能替代的。
  \item \textbf{「关节越轻越好控。」}——恰恰相反（\cref{sec:actrl-timescale}）。低惯量关节自然频率高、自然时间常数短，在常规控制率下\emph{欠采样发散}；对 $2\times10^{-4}$ 惯量腕施原始 PD，力矩冲到 $888\,\mathrm{N\!\cdot\!m}$。
  \item \textbf{「OSC 比关节空间控制高级，所以更好。」}——不一定。OSC 含 $\mathbf{M}^{-1}$，对 $\mathrm{cond}(\mathbf{M})$ 极敏感（\cref{sec:actrl-osc-illcond}）；异质惯量臂上力矩级 OSC 零空间收不住末端（漂 $60$--$110\,\mathrm{mm}$），而运动学级 resolved-rate 锁到 $0.00\,\mathrm{mm}$。先用条件数判可行性。
  \item \textbf{「零空间投影随便用个伪逆都行。」}——错（\cref{thm:actrl-null}）。只有\emph{动力学一致}逆 $\bar{\mathbf{J}}=\mathbf{M}^{-1}\mathbf{J}^\top\boldsymbol{\Lambda}$ 让零空间力矩对末端零加速度；用 Moore--Penrose $\mathbf{J}^+$ 会泄漏出末端寄生加速度。
  \item \textbf{「力矩臂要做柔顺必须装力传感器。」}——错（\cref{sec:actrl-imp-vs-adm}）。阻抗控制「测运动、出力」，靠主动施加弹簧律 $\boldsymbol{\tau}=\mathbf{J}^\top(-\mathbf{K}_c\Delta\mathbf{x})+\mathbf{g}$ 即呈现真阻抗，\emph{无需} F/T。
  \item \textbf{「$\mathbf{D}_c=2\sqrt{\mathbf{K}_c}$ 就是每个方向都临界阻尼。」}——只在「末端等效惯量=单位阵」近似下成立（\cref{sec:actrl-cartimp}）。精确临界阻尼需随操作空间惯量 $\boldsymbol{\Lambda}(\mathbf{q})$ 整形、是\emph{构型相关}的。
  \item \textbf{「混合力位可以沿世界系 $x,y,z$ 直接劈。」}——错（\cref{def:actrl-constraint-frame}）。必须在\emph{对齐接触几何}的约束坐标系里劈，否则力与位在同一轴纠缠。
\end{enumerate}

% ════════════════════════════════════════════════════════════════
\section{小结}
\label{sec:actrl-summary}

本章在上一部分控制理论地基（\cref{sec:ctrl-robot}）之上，把机械臂控制\emph{展开成完整工程谱系}：关节空间从 PD+重力补偿（调节，\cref{sec:actrl-pdg}）到计算力矩（跟踪，紧 $26\times$，\cref{sec:actrl-ctc}）；补进 P4 没有的异质惯量臂控制率选择（增益异质 + 解耦轻腕，\cref{sec:actrl-heterogeneous}）；任务空间给出 OSC 完整推导（动力学一致逆 + 零空间 + 对条件数敏感，\cref{sec:actrl-osc}）；接触控制涵盖阻抗/导纳（力矩臂真阻抗，\cref{sec:actrl-impedance}）与力位混合（选择矩阵 + 约束坐标系，\cref{sec:actrl-hybrid}）；最后给力矩跟踪的四签名诊断（\cref{sec:actrl-diagnostics}）。

\subsection*{符号表}
\begin{center}\small
\begin{longtable}{p{0.20\textwidth} p{0.50\textwidth} p{0.22\textwidth}}
\toprule
符号 & 含义 & 首现 \\
\midrule
\endhead
$\mathbf{q},\boldsymbol{\tau}\in\mathbb{R}^n$ & 关节角、关节力矩 & \cref{eq:actrl-manip} \\
$\mathbf{M},\mathbf{C},\mathbf{g}$ & 质量阵、科氏阵、重力力矩 & \cref{eq:actrl-manip} \\
$\mathbf{x}\in\mathbb{R}^m,\ \mathbf{J}$ & 末端任务坐标、任务雅可比 & \cref{def:actrl-spaces} \\
$\mathbf{e}=\mathbf{q}_d-\mathbf{q},\ \mathbf{a}_{\rm cmd}$ & 跟踪误差、计算力矩虚拟加速度 & \cref{eq:actrl-ctc} \\
$\mathbf{K}_p,\mathbf{K}_d$ & 比例、微分增益（可异质） & \cref{eq:actrl-ctc} \\
$\omega_n,\tau_{\rm nat}$ & 自然频率、自然时间常数 & \cref{eq:actrl-omegan} \\
$\boldsymbol{\Lambda}=(\mathbf{J}\mathbf{M}^{-1}\mathbf{J}^\top)^{-1}$ & 操作空间惯量矩阵 & \cref{eq:actrl-Lambda} \\
$\bar{\mathbf{J}}=\mathbf{M}^{-1}\mathbf{J}^\top\boldsymbol{\Lambda}$ & 动力学一致广义逆 & \cref{eq:actrl-jbar} \\
$\mathbf{N}=\mathbf{I}-\mathbf{J}^\top\bar{\mathbf{J}}^\top$ & 动力学一致零空间投影 & \cref{eq:actrl-N} \\
$\mathbf{F},\mathbf{F}^*$ & 末端广义力、力级 PD 反馈 & \cref{eq:actrl-osc-law} \\
$\mathbf{K}_c,\mathbf{D}_c$ & 笛卡尔刚度、阻尼 & \cref{eq:actrl-cartimp} \\
$\Delta\mathbf{x}=\mathbf{x}-\mathbf{x}_d$ & 末端位置偏差 & \cref{eq:actrl-cartimp} \\
$\mathbf{S},\mathbf{F}_d$ & 选择矩阵、期望接触力 & \cref{eq:actrl-hybrid} \\
$\mathrm{cond}(\mathbf{M})$ & 质量阵条件数 & \cref{sec:actrl-osc-illcond} \\
\bottomrule
\end{longtable}
\end{center}

\subsection*{定理与控制律速查}
\begin{center}\small
\begin{longtable}{p{0.30\textwidth} p{0.42\textwidth} p{0.22\textwidth}}
\toprule
名称 & 要点 & 标签 \\
\midrule
\endhead
PD+重力补偿 & 调节器，无源性 + LaSalle，仅需 $\mathbf{g}$ & \cref{thm:ctrl-pd-grav}（P4）\\
计算力矩 & $\boldsymbol{\tau}=\mathrm{RNEA}(\mathbf{q},\dot{\mathbf{q}},\mathbf{a}_{\rm cmd})$，闭环 $\ddot{\mathbf{e}}+\mathbf{K}_d\dot{\mathbf{e}}+\mathbf{K}_p\mathbf{e}=\mathbf{0}$ & \cref{eq:actrl-ctc}, \cref{thm:ctrl-ctc}（P4）\\
采样判据 & $\Delta t\gtrsim\tau_{\rm nat}$ 则欠采样发散 & \cref{thm:actrl-undersample} \\
腕耦合腐蚀 & 腕误差经 $M_{jk}$ 注入主关节，解耦 $a_{\rm cmd,\text{腕}}=0$ & \cref{thm:actrl-corruption} \\
操作空间动力学 & $\boldsymbol{\Lambda}\ddot{\mathbf{x}}+\boldsymbol{\mu}+\mathbf{p}=\mathbf{F}$ & \cref{eq:actrl-osc-dyn} \\
动力学一致零空间 & $\mathbf{N}\boldsymbol{\tau}_0$ 对末端零加速度 & \cref{thm:actrl-null} \\
笛卡尔阻抗 & $\boldsymbol{\tau}=\mathbf{J}^\top(-\mathbf{K}_c\Delta\mathbf{x}-\mathbf{D}_c\dot{\mathbf{x}})+\mathbf{g}$ & \cref{eq:actrl-cartimp} \\
力位混合 & $\mathbf{F}=\mathbf{S}\mathbf{F}_d+(\mathbf{I}-\mathbf{S})(\text{位 PD})$ & \cref{eq:actrl-hybrid} \\
\bottomrule
\end{longtable}
\end{center}

% ════════════════════════════════════════════════════════════════
\section{延伸阅读}
\label{sec:actrl-further}

\begin{itemize}
  \item \textbf{计算力矩与反馈线性化}：Spong 等\cite{spong2006robot} 给出标准证明与自适应扩展；反馈线性化的一般理论（相对阶、Lie 导数）见 Slotine--Li\cite{slotine1991applied} 与 \cref{sec:ctrl-robot-ctc} 的简述。
  \item \textbf{操作空间控制}：Khatib 1987 原文\cite{khatib1987operational} 是奠基；本书 \cref{ch:operational-space} 把它展开到浮动基、TSID 与稳定性证明，强烈推荐对照读。
  \item \textbf{阻抗控制}：Hogan 1985 三部曲\cite{hogan1985impedance} 奠定阻抗范式；笛卡尔阻抗的阻尼设计（构型相关临界阻尼）见 Albu-Schäffer 等\cite{albu2003cartesian}；无源性与能量端口的完整证明见 \cref{ch:force-wbc} 与 \cref{ch:operational-space}。
  \item \textbf{力位混合}：Raibert--Craig 1981 原文\cite{raibert1981hybrid} 提出选择矩阵；约束坐标系/自然--人工约束/接触稳定性的完整理论见 \cref{ch:force-wbc}（力控谱系）。
  \item \textbf{动力学工具}：质量阵 CRBA、逆动力学 RNEA、操作空间惯量算法见 Featherstone\cite{featherstone2008rbda} 与本部 \cref{ch:arm-dyn}。
\end{itemize}

% ════════════════════════════════════════════════════════════════
\section{后续关系}
\label{sec:actrl-next}

\begin{itemize}
  \item \textbf{向上承接}：本章的控制律建立在 \cref{ch:control} 的控制理论、\cref{ch:arm-kin} 的雅可比与奇异、\cref{ch:arm-dyn} 的质量阵与条件数之上；轨迹由 \cref{ch:arm-traj} 提供（计算力矩跟踪的 $\mathbf{q}_d(t)$、力矩约束时间参数化）。
  \item \textbf{落到实战}：本章全部控制律在 \cref{ch:arm-prac} 的六轴力矩臂 arm\_dm 上跑通——重力补偿、笛卡尔阻抗、计算力矩、混合力位、运动学 vs 力矩级零空间对照、病态根因，并组装成「规划 $\to$ 计算力矩 $\to$ 阻抗」的全栈连续会话。
  \item \textbf{启下到强化学习}：经典控制（OSC、阻抗、计算力矩）是 \cref{ch:rlmc-manip} 强化学习操作的\emph{地基}——RL 操作里「OSC 动作层」「隐式学到的阻抗」「DiffIK 的阻尼伪逆」都假定了本章与 \cref{ch:arm-kin,ch:arm-dyn} 已建立的经典理论。读者可把本章当作进入 RL 操作前的「经典对照系」。
\end{itemize}

% ════════════════════════════════════════════════════════════════
\section{故障排查}
\label{sec:actrl-troubleshoot}

\begin{table}[H]\centering\small
\caption{机械臂控制常见故障速查}
\label{tab:actrl-troubleshoot}
\begin{tabularx}{\linewidth}{p{0.26\textwidth} L L}
\toprule
现象 & 根因 & 对策 \\
\midrule
跟踪有周期滞后 & PD+重力补偿缺前馈 & 换计算力矩，补 $\mathbf{M}\ddot{\mathbf{q}}_d+\mathbf{C}\dot{\mathbf{q}}_d$（\cref{eq:actrl-ctc}）\\
轻腕高频发散 & 自然频率 $\gg$ 控制率，欠采样 & 降腕增益（异质 $\mathbf{K}_p$）+ 解耦 $a_{\rm cmd,\text{腕}}=0$（\cref{algo:actrl-hetero-ct}）\\
主关节力矩异常大、卡住 & 腕误差经 $M_{jk}$ 耦合腐蚀 & 解耦腕（\cref{thm:actrl-corruption}）\\
OSC 末端收不住、漂 60--110\,mm & 质量阵病态，投影性被破坏 & 改运动学级 resolved-rate（无 $\mathbf{M}^{-1}$，\cref{sec:actrl-osc-illcond}）\\
零空间次任务推动末端 & 用了 Moore--Penrose 伪逆 & 改动力学一致逆 $\bar{\mathbf{J}}$（\cref{def:actrl-jbar}）\\
阻抗末端超调振荡 & 阻尼不足/非临界 & $\mathbf{D}_c=2\sqrt{\mathbf{K}_c}$，必要时随 $\boldsymbol{\Lambda}$ 整形（\cref{eq:actrl-critdamp}）\\
混合控制趋近时飞/卡 & 关节切分 + 折叠位子雅可比奇异 & 全雅可比 + 关节 CT 展开 + 选 reach-free 扫掠轴（\cref{sec:actrl-selection}）\\
home 位不悬浮、关节下垂 & 限位夹持（非重力补偿错）& 数据驱动诊断：零力矩落体 + 已知力矩权限（\cref{sec:actrl-pdg}）\\
\bottomrule
\end{tabularx}
\end{table}

\begin{practice}[练习]
\begin{enumerate}
  \item （承接 P4）复用 \cref{thm:ctrl-ctc}，把计算力矩律 \cref{eq:actrl-ctc} 代入 \cref{eq:actrl-manip}，验证模型精确时闭环误差满足 $\ddot{\mathbf{e}}+\mathbf{K}_d\dot{\mathbf{e}}+\mathbf{K}_p\mathbf{e}=\mathbf{0}$；再讨论：若 $\hat{\mathbf{M}}=\mathbf{M}+\Delta\mathbf{M}$ 有偏，残差项是什么形式？
  \item （异质惯量）给定腕惯量 $I=2\times10^{-4}$、阻尼 $b=0.1$，估其自然时间常数 $\tau_{\rm nat}$ 与对应频率；若控制率为 $200\,\mathrm{Hz}$，判断是否欠采样（用 \cref{thm:actrl-undersample}）。再算：施 $\mathbf{K}_p=150$ 的 PD 时 $\omega_n=\sqrt{k/I}$ 为多少？
  \item （OSC 推导）从 \cref{eq:actrl-osc-xddot} 与关节动力学出发，自行补全 \cref{eq:actrl-osc-dyn} 中 $\boldsymbol{\mu},\mathbf{p}$ 的表达式，并验证 $\bar{\mathbf{J}}$ 满足 $\mathbf{J}\bar{\mathbf{J}}=\mathbf{I}$。
  \item （零空间）证明 \cref{eq:actrl-N} 的 $\mathbf{N}$ 是\emph{投影}（$\mathbf{N}^2=\mathbf{N}$），并说明为何正则化 $\mathbf{M}+\lambda\mathbf{I}$ 会破坏这个性质。
  \item （阻抗实证）阻抗刚度 $\mathbf{K}_c=250\,\mathrm{N/m}$，施外力 $50\,\mathrm{N}$，按线性弹簧律预测稳态退让，并与实测 $199.9\,\mathrm{mm}$ 比较。若要退让减半，$\mathbf{K}_c$ 应取多少？
  \item （力位混合）一个孔轴沿世界系 $(1,1,0)/\sqrt{2}$ 方向的插装任务，写出约束坐标系的选取、自然约束与人工约束（\cref{def:actrl-constraint-frame}），并给出选择矩阵 $\mathbf{S}$。
  \item （故障诊断）观察到某关节力矩朝一个方向单调增大直至饱和、关节越跑越快。对照 \cref{tab:actrl-signatures}，这是哪种签名？列两个最可能的根因与各自的验证手段。
\end{enumerate}
\end{practice}
